\documentclass[compacto,10pt]{pubciencias-notas}

% -- Paquetes adicionales
\usepackage{comandosEPN}

% -- Datos del libro
\autor{Departamento de Formación Básica}
\materia{Álgebra Lineal}
\nota{Semana 02 Clase 01}
\fecha{Semestre 2025-A}

\logouno[1.1cm]{Logo_EPN.eps}
\logodos[1.7cm]{Logo_DFB.eps}

\definecolor{colordef}{cmyk}{0.81,0.62,0.00,0.22}

\begin{document}

\encabezado

\vspace{-8mm}

\section{Matrices}
\setcounter{subsection}{2}
\subsection{Tipos de matrices}

\begin{defi}[Matriz triangular]
    Sea $A\in \mathbb{K}^{n \times n}$. Se dice que $A$ es
    \begin{itemize}
    \item 
        una matriz triangular superior si $a_{ij} = 0 $ para todo $i>j$ con $i,j \in J$; y
    \item 
        una matriz triangular inferior si $a_{ij} = 0 $ para todo $i<j$ con $i,j \in J$.
    \end{itemize}
\end{defi}

\begin{ejem}
A continuación, un ejemplo de una matriz triangular superior
\[
	\begin{pmatrix}3&2\\0&1\end{pmatrix}
\]
y de una matriz triangular inferior:
\[
	\begin{pmatrix}1&0&0\\-8&2&0\\0&-6&3\end{pmatrix}.
\]
\end{ejem}

\begin{defi}[Matriz diagonal]
    Sea $A\in \mathbb{K}^{n \times n}$. Se dice que $A$ es una matriz diagonal si verifica que:
    \[ 
        a_{ij} = 0
    \]
    para todo $i\in J$ y $j\in J$ con $i\neq j$.
\end{defi}

\begin{ejem}
Vemos algunos ejemplos de matrices diagonales:
\[
	\begin{pmatrix}
	-2&0\\
	0&5
	\end{pmatrix},
	\quad
 	\begin{pmatrix}
	1&0&0\\
	0&1&0\\
	0&0&1
	\end{pmatrix},
	\quad
	\begin{pmatrix}
	-3&0&0&0\\
	0&-3&0&0\\
	0&0&3&0\\
	0&0&0&3\\
	\end{pmatrix}.
\]
Notemos que en todos los casos los términos fuera de la diagonal principal son todos iguales a cero. 
\end{ejem}

\begin{comment}
\begin{defi}[Matriz escalar]
    Sea $A\in \mathbb{K}^{n \times n}$. Se dice que $A$ es una matriz escalar si es una matriz diagonal que verifica:
    \[ 
        a_{ii} = \alpha
    \]
    para todo $i\in J$, con $\alpha\in\K$, es decir,  todos los elementos de la diagonal son iguales.
\end{defi}

\begin{ejem}
Vemos algunos ejemplos de matrices diagonales:
\[
	\begin{pmatrix}
	-2&0\\
	0&-2
	\end{pmatrix},
	\quad
	\begin{pmatrix}
	5&0&0&0\\
	0&5&0&0\\
	0&0&5&0\\
	0&0&0&5\\
	\end{pmatrix}.
 \]
\end{ejem}

\begin{advertencia}
Note que una matriz identidad es un tipo de matriz escalar. Además, la matriz escalar se puede expresar como la multiplicación de un escalar por la matriz identidad.
\end{advertencia}

\begin{ejem} Considere a la siguiente matriz diagonal
\[
	A= \begin{pmatrix}
	3&0\\
	0&3\\
	\end{pmatrix}
	=3 \begin{pmatrix}
	1&0\\
	0&1\\
	\end{pmatrix}
	=3 I_2.
\]
Note que la matriz escalar dada se escribe como el producto del escalar 3 con la matriz identidad.
\end{ejem}

\begin{ejer}\hspace{0pt} Demuestre que la suma de dos matrices escalares es una matriz escalar. ¿Qué se puede decir sobre la resta?
\end{ejer}
\end{comment}

\begin{defi}[Matriz simétrica]
    Sea $A\in \mathbb{K}^{n \times n}$. Se dice que $A$ es simétrica si $A^\intercal = A$.
\end{defi}

\begin{ejem}\hspace{0pt} La matriz 
\[
A = \begin{pmatrix}0&1&2&2\\1&-8&2&3\\2&2&5&4\\2&3&4&1
\end{pmatrix},
\]
es simétrica, en efecto, tenemos que
\[
    A^{\intercal} = A.
\]
\end{ejem}

\begin{ejer}\hspace{0pt} Dada la matriz $A\in\Mat[\K]{n}{n}$, muestre que $A^{\intercal} + A$ es simétrica.
\end{ejer}

\begin{defi}[Matriz antisimétrica]
    Sea $A\in \mathbb{K}^{n \times n}$. Se dice que $A$ es antisimétrica si $A^\intercal = -A$.
\end{defi}

\begin{ejem}\hspace{0pt} La matriz
\[
A = \begin{pmatrix}0&1&-8\\-1&0&3\\8&-3&0\end{pmatrix},
\]
es antisimétrica, en efecto, tenemos que
\[
    A^{\intercal} = - A.
\]
\end{ejem}

\begin{ejer}\hspace{0pt} Dada la matriz $A\in\Mat[\K]{n}{n}$, muestre que $A^{\intercal} - A$ es antisimétrica.
\end{ejer}

\begin{comment}
\begin{defi}[Matriz nilpotente]
    Sea $A\in \mathbb{K}^{n \times n}$ y $r \in \N$. Se dice que $A$ es nilpotente de orden $r$ si $r$ es el menor entero positivo tal que
    \[ 
        A^r = 0.
    \]
\end{defi}

\begin{ejem}\hspace{0pt} La siguiente matriz es nilpotente de orden 2:
\[
	A=\begin{pmatrix}1&1\\-1&-1\end{pmatrix}
\]
En efecto, determinemos a $A^2$,
\[
	A^2=A\cdot A=
	\begin{pmatrix}1&1\\-1&-1\end{pmatrix}\cdot 
    \begin{pmatrix}1&1\\-1&-1\end{pmatrix} =
    \begin{pmatrix}0&0\\0&0\end{pmatrix}.
\]
Luego, la matriz $A$ es nilpotente de orden $2$.
\end{ejem}

\begin{defi}[Matriz idempotente]
    Sea $A\in \mathbb{K}^{n \times n}$. Se dice que $A$ es una matriz idempotente si $A^2 = A$.
\end{defi}

\begin{ejem}\hspace{0pt} La matriz 
\[
	A=\begin{pmatrix}1&0\\3&0\end{pmatrix}
\]
es idempotente, en efecto, calculemos a $A^2$,
\[
	A^2=
	\begin{pmatrix}1&0\\3&0\end{pmatrix}\cdot \begin{pmatrix}1&0\\3&0\end{pmatrix}=
	\begin{pmatrix}1&0\\3&0\end{pmatrix}=A.
\]
\end{ejem}

\begin{advertencia}
Note que si la matriz $A$ es idempotente entonces se puede concluir que $A^n=A$, para todo $n \in \Nbb^*$.
\end{advertencia}
\end{comment}

\begin{defi}[Matriz escalonada por filas]
Se dice que una matriz está en forma escalonada por filas cuando satisface las siguientes propiedades:
\begin{enumerate}
\item Todas las filas que constan de ceros, si las hay, están en la parte inferior de la matriz. 
\item La entrada principal (primer elemento distinto de cero) de cada fila está en una columna a la derecha de la entrada principal de la fila de arriba. 
\item En una columna, todas las entradas abajo de una entrada principal son ceros.
\end{enumerate}
\end{defi}

\begin{defi}[Matriz escalonada reducida por filas]
Se dice que una matriz está en forma escalonada reducida por filas cuando satisface las siguientes propiedades:
\begin{enumerate}
\item Es una matriz escalonada por filas. 
\item La entrada principal de cada fila es 1.
\item Cada 1 principal es el único elemento distinto de cero de su columna
\end{enumerate}
\end{defi}

\begin{ejem}\hspace{0pt} Las siguientes matrices son escalonadas por filas
\[
    \begin{pmatrix}
        1&4\\
        0&1\\
        0&0
    \end{pmatrix},\qquad
    \begin{pmatrix}
        1&6\\
        0&2
    \end{pmatrix} \texty
    \begin{pmatrix}
        3&0&1&-2&-1\\
        0&0&1&4&-2\\
        0&0&0&-7&2
    \end{pmatrix}.
\]

\end{ejem}
\begin{ejem}\hspace{0pt} Las siguientes matrices son escalonadas reducidas por filas
\[
    \begin{pmatrix}
        1&0&0&2\\
        0&1&0&-2\\
        0&0&1&0\\
        0&0&0&0
    \end{pmatrix},\qquad
    \begin{pmatrix}
        1&6\\
        0&0
    \end{pmatrix} \texty
    \begin{pmatrix}
        1&0&0&0&-1\\
        0&0&1&0&-2\\
        0&0&0&1&2
    \end{pmatrix}.
\]
\end{ejem}

\subsection{Equivalencia por filas}

\begin{defi}[Operaciones elementales de fila]
    Dada una matriz $A\in\K^{m \times n}$, una operación elemental por filas sobre $A$ es una de las siguientes:
    \begin{itemize}
    \item   
        \textbf{Intercambio de filas:} dados $i , j \in I$, intercambiar la fila $i$ por la fila $j$, denotado por
        \[
            F_i \leftrightarrow F_j,
        \]
        es reemplazar la fila
        \[
            \begin{pmatrix}a_{i1}&a_{i2}&\ldots& a_{in}\end{pmatrix}
        \]
        por la fila
        \[
            \begin{pmatrix}a_{j1}&a_{j2}&\ldots& a_{jn}\end{pmatrix}
        \]
        y viceversa.
        \item 
        \textbf{Multiplicar una fila por un escalar:} dados $i\in I$ y $\alpha\in \K$, con $\alpha\neq 0$ y $\alpha \neq 1$, multiplicar la fila $i$ por $\alpha$, denotado por
        \[
            \alpha F_i \rightarrow F_i,
        \]
        es reemplazar la fila
        \[
            \begin{pmatrix}a_{i1}&a_{i2}&\ldots& a_{in}\end{pmatrix}
        \]
        por 
        \[
            \begin{pmatrix}\alpha a_{i1}&\alpha a_{i2}&\ldots&\alpha a_{in}\end{pmatrix}.
        \]
       \item 
        \textbf{Sumar un múltiplo de una fila con otra:} dados $i,j\in I$ y $\alpha\in \K$, multiplicar la fila $i$ por $\alpha$ y sumarlo a la fila $j$, denotado por
        \[
            \alpha F_i + F_j\rightarrow F_j,
        \]
        es reemplazar la fila
        \[
            \begin{pmatrix}a_{j1}&a_{j2}&\ldots& a_{jn}\end{pmatrix}
        \]
        por 
        \[
            \begin{pmatrix}\alpha a_{i1} + a_{j1}&\alpha a_{i2} + a_{j2}&\ldots&\alpha a_{in} + a_{jn}\end{pmatrix}.
        \]
        
\end{itemize}
\end{defi}

\begin{ejem}\hspace{0pt}\label{Ejem:OprEle} \emph Dada la matriz:
     \[
           A=\begin{pmatrix}
           1&2&3 \\
           0&-1&2\\
           -1&2&-3
           \end{pmatrix}
        \]
    Aplicar las siguientes operaciones
    \begin{itemize}
        \item \textit{Intercambio de filas:} $F_2 \leftrightarrow F_3$

        Se obtiene la siguiente matriz
        \[
           B=\begin{pmatrix}
           1&2&3 \\
           -1&2&-3\\
           0&-1&2
           \end{pmatrix}
        \]

        \item \textit{Multiplicar una fila por un escalar:} $3  F_1 \rightarrow F_1$

        Se obtiene la siguiente matriz
        \[
           C=\begin{pmatrix}
           3&6&9 \\
           0&-1&2\\
           -1&2&-3
           \end{pmatrix}
        \]

        \item \textit{Sumar un múltiplo de una fila con otra:} $3 F_1 + F_3\rightarrow F_3$

        Se obtiene la siguiente matriz
        \[
           D=\begin{pmatrix}
           1&2&3 \\
           0&-1&2\\
           2&8&6
           \end{pmatrix}
        \]
    \end{itemize}
    \end{ejem}
    
% \begin{ejem}\hspace{0pt} \emph{Multiplicar una fila por un escalar.} Considerando a la matriz $A$ del ejemplo \ref{Ejem:OprEle}, si se aplica la operación de fila $3  F_1 \rightarrow F_1$ se tiene la siguiente matriz:
%     \[
%            C=\begin{pmatrix}
%            3&6&9 \\
%            0&-1&2\\
%            -1&2&-3
%            \end{pmatrix}
%         \]
    
%     \end{ejem}
    
% \begin{ejem}\hspace{0pt}\label{Ejem:OpeEle1} \emph{Sumar un múltiplo de una fila con otra.} Considere a la matriz $A$ del ejemplo \ref{Ejem:OprEle}, si se aplica la operación de fila $3 F_1 + F_3\rightarrow F_3$ se tiene la siguiente matriz:
%     \[
%            D=\begin{pmatrix}
%            1&2&3 \\
%            0&-1&2\\
%            2&8&6
%            \end{pmatrix}
%         \]
    
% \end{ejem}


\begin{ejer}\hspace{0pt}\label{Ejer:OpeELe}   
Dada la matriz
\[
    A=\begin{pmatrix}
           1&2&3 \\
           0&-1&2\\
           -1&2&-3
           \end{pmatrix}.
\]
Note que en el ejemplo anterior se realizaron tres operaciones sobre esta matriz. Ahora, realizar las mismas operaciones sobre $A$ de manera consecutiva. 
\begin{enumerate}
\item $F_{2} \leftrightarrow F_{3}$.
\item $3F_{1} \rightarrow F_{1}$.
\item $3F_{1} + F_{3} \rightarrow F_{3}$.
\end{enumerate}
¿Qué matriz se ha obtenido?
\end{ejer}

% \begin{teo}
% Dada una matriz $A \in \Mat[\K]{m}{n}$, a cada operación elemental de filas $e$, le corresponde una operación elemental de filas $e^{- 1}$, del mismo tipo de $e$, tal que, $e(e^{- 1}(A))=e^{- 1}(e(A))=A$.

% \end{teo}
% \begin{advertencia}
%     Notemos que si $e_1, e_2, e_3$ son operaciones de fila tal que:
% \[
% e_1:F_i \leftrightarrow F_j,
% \]
% \[
% e_2:\alpha \cdot F_i \rightarrow F_i,
% \] 
% \text{y}
% \[
% e_3:\alpha \cdot F_i +F_j\rightarrow F_j.
% \]

% Entonces existen las operaciones elementales de fila  $e_1^{-1}$,$e_2^{-1}$,$e_3^{-1}$ donde:
% \[
% e_1^{-1}:F_i \leftrightarrow F_j,
% \]
% \[
% e_2^{-1}:\dfrac{1}{\alpha} \cdot F_i \rightarrow F_i,
% \] 
% \text{y}
% \[
% e_3^{-1}:-\alpha \cdot F_i +F_j\rightarrow F_j.
% \]
% \end{advertencia} 
% \begin{advertencia}
% A la operación elemental $e^{- 1}$ se le llama operación elemental inversa de $e$.
% \end{advertencia}

\begin{defi}[Equivalente por filas]
    Sean $A,B\in \Mat[\K]{m}{n}$, se dice que la matriz $A$ es equivalente por filas a una matriz $B$, denotado por $A\sim B$, si se puede obtener $B$ al aplicar a la matriz $A$ una sucesión de operaciones elementales de fila.
\end{defi}

\begin{advertencia}
Si $A$ es equivalente por filas a $B$, entonces es $B$ equivalente por filas a $A$.
\end{advertencia}

\begin{ejem}\hspace{0pt} Al aplicar la operación $2F_2+F_1\rightarrow F_1$ a la matriz
\[
    A=\begin{pmatrix}
        2&3&1&1\\
        1&-2&-1&3\\
        0&2&-3&4
    \end{pmatrix}
\]
se obtiene la matriz
\[
    B=\begin{pmatrix}
        4&-1&-1&7\\
        1&-2&-1&3\\
        0&2&-3&4
    \end{pmatrix}.
\]
Entonces $A$ es equivalente por filas a $B$, es decir, $A\sim B$. 
\end{ejem}

% \begin{ejer}\hspace{0pt} Considere a la matriz $A$ y a las operaciones elementales de fila del ejercicio \ref{Ejer:OpeELe}, escriba a las correspondientes operaciones elementales de fila inversas.
% \end{ejer}

\end{document}
